% --- Part 6: Proposed GPU Library (Slides 22-26) ---
\section{Propuesta de Trabajo}

\begin{frame}{Limitaciones de las Herramientas Actuales}
    \begin{itemize}
        \item Bibliotecas como SEAL, HELib o PALISADE usan solo CPU.
        \item El cómputo homomórfico es intensivo; la CPU es un cuello de botella.
        \item Necesidad de aceleración por hardware para modelos complejos.
    \end{itemize}
\end{frame}

\begin{frame}{Stack Tecnológico}
    \begin{itemize}
        \item \textbf{FIDESlib:} Backend optimizado en CUDA/C++.
        \item \textbf{OpenFHE-Python:} Interfaz de alto nivel.
        \item \textbf{Pybind11:} Bindings para exponer C++ en Python.
    \end{itemize}
\end{frame}

\begin{frame}{Objetivos del Trabajo}
    \begin{itemize}
        \item Desarrollar una biblioteca Python para tensores cifrados en GPU.
        \item Abstraer la complejidad criptográfica del usuario.
        \item Habilitar inferencia privada de modelos reales.
    \end{itemize}
\end{frame}

\begin{frame}{Casos de Uso y Benchmarking}
    \begin{itemize}
        \item Implementación de redes DNN y CNN cifradas.
        \item Comparar contra TenSEAL (CPU) y Concrete ML (GPU-TFHE).
        \item Evaluar impacto de funciones de activación y parámetros.
    \end{itemize}
\end{frame}

\begin{frame}{Resumen y Trabajo Futuro}
    \begin{itemize}
        \item Reducción de la brecha entre criptografía y Machine Learning.
        \item Provisión de herramientas para computación privada eficiente.
        \item Extensión a modelos iterativos y procesamiento en tiempo real.
    \end{itemize}
\end{frame}
